\newpage
\section{Formelsammlung zur Physik des Kontinuums}\label{kap:kontmech_formelsammlung}
Im Folgenden geben wir zum Arbeitsbuch einige h�ufig benutzte Formeln und Beziehungen der Himmelsmechanik.
\subsection{Vektoranalysis}
Es ist vorteilhaft, Differentialoperatoren nicht nur in kartesischen, sondern auch, je nach Geometrie, in Zylinder- und Kugelkoordinaten anzugeben. \\

\textbf{Ableiten von skalaren Feldern} $\phi\lk\vec{r}\rk$ \textbf{und Vektorfeldern} $\vec{A}\lk \vec{r}\rk$
\begin{align*}
    \text{Gradient:}  ~~ &\gradv f = \nablav f ~~~~ \text{\bzw} ~~~~ \lk \nablav f\rk_i = \d_i f \quad \eqref{eq:kontmech_gradient}\\
    \text{Divergenz:} ~~ &\divv \cdot \vec{A} = \nablav\cdot\vec{A} = \d_i\,F_i \quad \eqref{eq:kontmech_divergenz}\\
    \text{Rotation:}  ~~ &\rotv\,\vec{A} = \nablav\times \vec{A} ~~~~ \text{bzw.} ~~~~ \lk \rotv \vec{A}\rk_i = \epsilon_{ijk}\,\d_j\,F_k \quad \eqref{eq:kontmech_rotation}
\end{align*}
\vspace{0.5cm}
\textbf{Zweite Ableitungen von skalaren Feldern} $ f\lk r\rk$ \textbf{und Vektorfeldern} $\vec{A}\lk r\rk$
\begin{align*}
    \divv\,\gradv f &= \nablav\,\nablav f = \Delta  f  \quad \eqref{eq:kontmech_nablaop05} \\
    \gradv\,\divv\,\vec{A} &= \nablav\,\lk \nablav\,\vec{A}\rk \\
    \divv \,\rotv\,\vec{A} &= \nablav\,\lk \nablav\times \vec{A}\rk = 0 \quad \eqref{eq:kontmech_nablaop06} \\
    \rotv \gradv\, f &= \nablav\times\lk \nablav f \rk = 0 \quad \eqref{eq:kontmech_nablaop07} \\
    \rotv\,\rotv\,\vec{A} &= \nablav\times\lk \nablav\times\vec{A}\rk = \nablav\,\lk \nablav\,\vec{A}\rk - \Delta \vec{A} \quad \eqref{eq:kontmech_nablaop04} \\
\end{align*}
\vspace{0.5cm}
\textbf{Differentialoperatoren in Zylinderkoordinaten}
\begin{align*}
    \nablav  f &= \vec{e}_\rho \,\frac{\d f}{\d \rho} + \vec{e}_\phi \,\frac{1}{\rho}\frac{\d f}{\d \phi} + \vec{e}_z\frac{\d f}{\d z} \\
    \Delta f &= \underbrace{\frac{1}{\rho}\,\frac{\d}{\d \rho} \lk \rho\,\frac{\d f}{\d \rho }\rk}_{\frac{\d^2 f}{\d \rho^2} + \frac{1}{\rho}\,\frac{\d f}{\d \rho}}+ \frac{1}{\rho^2}\,\frac{\rho^2 f}{\d \phi^2} + \frac{\d^2 f}{\d z^2}
\end{align*}
\vspace{0.5cm}
\textbf{Differentialoperatoren in Kugelkoordinaten}
\begin{align*}
    \nablav f &= \vec{e}_r \,\frac{\d f}{\d r} + \vec{e}_\theta \,\frac{1}{r} \frac{\d f}{\d \phi} + \vec{e}_\phi \,\frac{1}{r\sin\theta}\frac{\d f}{\d \phi} \\
    \Delta f &= \underbrace{\frac{1}{r^2}\,\frac{\d}{\d r}\lk r^2\,\frac{\d f}{\d r} \rk}_{\frac{\d^2 f}{\d r^2} + \frac{2}{r}\,\frac{\d f}{\d r}}
		+ \frac{1}{r^2\sin\theta} \frac{\d }{\d \theta} \lk \sin\theta \frac{\d f}{\d \theta} \rk + \frac{1}{r^2 \sin^2\theta}\frac{\d^2 f}{\d \phi^2}
\end{align*}
\vspace{0.5cm}

\newpage
\subsection{Wichtige Gleichungen der Kontinuumsmechanik}

\textbf{Kontinuit�tsgleichung}
\begin{equation*}
  \frac{\partial}{\partial t}\varrho + \divv \vec{j} = \frac{\partial}{\partial t}\varrho + \divv \lk \varrho\, \vec{v} \rk = 0 \qquad  \eqref{eq:kontmech_KontinuitaetsglDiff}
\end{equation*}
\vspace{0.5cm}
\textbf{Euler-Gleichung der Kontinuumsmechanik}
\begin{equation*}
  \varrho \lk \lk \vec{v} \cdot \nabla \rk \vec{v} + \frac{\partial}{\partial t}
  \vec{v} \rk = \vec{f}_\text{V}  + \vec{f}_\text{R} - \vec{\nabla} p \; \qquad
  \eqref{eq:kontmech_EulerGleichung}
\end{equation*}
\vspace{0.5cm}
\textbf{Bernoulli-Gleichung}
\begin{equation*}
  \frac{1}{2} \varrho \, u_1^2 + p_1 = \frac{1}{2} \varrho \, u_2^2 +  p_2  \qquad  \eqref{eq:kontmech_BernoulliGl}
\end{equation*}
\vspace{0.5cm}
\textbf{Navier-Stokes-Gleichung}
\begin{equation*}
  \varrho \lk \lk \vec{v} \cdot \nablav \rk \vec{v} + \frac{\partial}{\partial t}
  \vec{v} \rk = \vec{f}_\text{V}  + \eta \, \Delta \vec{v} - \nablav p \; 
  \qquad \eqref{eq:kontmech_NavierStokes}
\end{equation*}
\vspace{0.5cm}
\textbf{Navier-Stokes-Gleichung - skaliert}
\begin{equation*}
  \lk \tilde{\vec{v}} \cdot \tilde{\nablav} \rk \tilde{\vec{v}}
  + \frac{\partial}{\partial \tilde{t}} \tilde{\vec{v}} 
  = \frac{1}{\mathrm{Re}} \, \tilde{\Delta} \tilde{\vec{v}}
  - \tilde{\nablav} \tilde{p} \; \qquad \eqref{eq:kontmech_NavierStokesSkaliert02}
\end{equation*}
\vspace{0.5cm}
\textbf{Reynoldszahl}
\begin{equation*}
  \mathrm{Re} = \frac{\varrho\,u\,l}{\eta}
\end{equation*}
\vspace{0.5cm}

\newpage
\subsection{Wichtige Gleichungen der Wellentheorie}

\textbf{Wellengleichung f�r die schwingende Saite}
\begin{equation*}
  \varrho \frac{\partial^2 y(x,t)}{\partial t^2}
  - E \frac{\partial^2 y(x,t)}{\partial x^2} = 0 \; \qquad \eqref{eq:kontmech_PDGLSaite2}
\end{equation*}
\vspace{0.5cm}
\textbf{Gruppengeschwindigkeit}
\begin{equation*}
  v_\text{Gr} = \frac{\D \omega(k)}{\D k} \qquad \eqref{eq:kontmech_DispersionsRel04}
\end{equation*}
\vspace{0.5cm}

\textbf{Korteweg-de-Vries-Gleichung}
\begin{equation*}
  \frac{\partial \psi}{\partial t} + 6\,\psi\,\frac{\partial \psi}{\partial x} + \frac{\partial^3 \psi}{\partial x^3} = 0\,\qquad \eqref{eq:kontmech_soliton01}
\end{equation*}
\vspace{0.5cm}
\textbf{Sinus-Gordon-Gleichung}
\begin{equation*}
  \frac{1}{c^2}\frac{\partial^2 \psi}{\partial t^2}  - \frac{\partial^2 \psi}{\partial x^2}  + \sin \psi = 0\, \qquad \eqref{eq:kontmech_soliton03}
\end{equation*}
\vspace{0.5cm}





